\documentclass[11pt]{article}
\usepackage{amsmath,amssymb,amsthm}
\usepackage[margin=1in]{geometry}
\usepackage{amsmath,amssymb,amsthm,mathtools}
\usepackage{enumitem}
\usepackage[pdfauthor={},pdftitle={Solution},pdfcreator={},pdfproducer={}]{hyperref}

\newtheorem{theorem}{Theorem}
\newtheorem{proposition}{Proposition}[section]
\newtheorem{lemma}[proposition]{Lemma}
\newtheorem{corollary}[proposition]{Corollary}
\theoremstyle{definition}
\newtheorem{definition}[proposition]{Definition}
\newtheorem{remark}[proposition]{Remark}

\newcommand{\Teich}{\mathrm{Teich}}
\newcommand{\PMod}{\mathrm{PMod}}
\newcommand{\Ch}{\hat{\mathbb C}}
\newcommand{\Rh}{\hat{\mathbb R}}
\newcommand{\HH}{\mathbb H}
\newcommand{\D}{\mathcal D}
\newcommand{\M}{\mathcal M}
\newcommand{\Push}{\mathrm{Push}}
\newcommand{\Img}{\operatorname{Im}}
\newcommand{\Rea}{\operatorname{Re}}

\newcommand{\fptitle}{Review of Problem 7 -- Siege}
\newcommand{\reviewtag}{Reviewer 3}
\usepackage{fp-review}

\begin{document}
\maketitle

\section{Review}

Due to:
\begin{itemize}
    \item The extremely poor choice of notation;

    \item Abundance of overly detailed numerics and inline computations;

    \item Poor structuring of the arguments within a given proof or between multiple statements;
\end{itemize}
It is unreasonably hard to ascertain the correctness of the proposed solution. It is possible that it internally correct; the initial strategy for determining the criteria that the curve has to satisfy in order to compose a proper solution follows the general themes of the relevant papers (Hironaka, Koch) and (Firsova, Kahn, Selinger). It also seems to correctly arrive at the first part of the problem, the preservation of the deformation space under the Dehn twist. 

The extra steps necessary, that is, explicitly constructing the curve which permutes the components along an infinite orbit, understanding the lift of its Dehn twist to a deck transformation on the Hurwitz space (corresponding to some loop within it), and the computation of the winding number of said loop on an appropriate space, are extremly difficult to parse through given the aforementioned characteristics of the submission. The submission also spends an unnecessary amount of rederiving well known facts within the field that can just be properly attributed and referenced, and seems to make mistakes in identifying points in Teichmuller space with their coordinates in Moduli space. Yet another problem is missing attributions, and possibly irrelevant citations corresponding to the poor choice of notational strategy, as remarked on the inline reviews.

\paragraph{Correctness.}
\begin{itemize}
	\item  1: essentially flawless mathematical argument
	\item  0.5: minor errors that are easily corrected
	\item  0: blank solution, or acknowledged partial result
	\item  $-0.5$: major errors
	\item  $-1$: proposed proof makes no progress towards the result.
\end{itemize}
Note that points should be docked for wrong statements/ solutions.

\noindent Score: N/A (correctness could not be determined)

\paragraph{Exposition.}
\begin{itemize}
	\item  1: well-written and clear solution.
	\item  0: argument that is possible to follow with serious effort
	\item  $-1$: very poor exposition
\end{itemize}

\noindent Score: $-1$

\paragraph{Attributions.}
\begin{itemize}
	\item  1: solutions contains fully appropriate attributions
		to related work and theorems being used.
	\item  0: moderate issues with missing attributions or hallucinated references
	\item  $-1$: blatant plagiarism or hallucinations
\end{itemize}

\noindent Score: 0

\newpage
\section{Submission with inline comments}


\section{The problem}

Let $f\colon S^2\to S^2$ be a branched cover, $V(f)$ the set of its critical values, and $A,B\subset S^2$ finite sets.
\begin{enumerate}
\item If $f(A)\cup V(f)\subset B$, then every complex structure on $(S^2,B)$ can be pulled back to a complex structure on $(S^2,A)$, and this induces the map $f^*\colon\Teich(B)\to\Teich(A)$.
\item If $A\subset B$, then forgetting the points $B\setminus A$ induces the map $i_*\colon\Teich(B)\to\Teich(A)$.
\end{enumerate}
If $f(A)\cup V(f)\cup A\subset B$, then the \emph{deformation space} is defined by
\[\D_f:=\{\tau\in\Teich(B): f^*(\tau)=i_*(\tau)\}.\]
Consider a quadratic rational map $f$ with critical points $c_1$ and $c_2$. Assume that $c_1$ belongs to a cycle of length $4$, and the critical value $v=f(c_2)$ does not belong to the cycle. By taking
\[A=\{c_1,f(c_1),f^{\circ2}(c_1),f^{\circ3}(c_1)\},\qquad B=A\cup\{v\},\]
we get a deformation space $\D_f\subset\Teich(B)$. Its image in the moduli space is the set of conjugacy classes of quadratic rational maps with the properties described above.

\emph{Explicitly describe a simple closed curve $\gamma\subset(S^2,B)$ such that a power of the Dehn twist around this curve preserves the deformation space, and its action on the set of connected components of $\D_f$ has an infinite orbit.}

\section{The result}

Throughout, $f$ is as in the problem, $a_i=f^{i}(c_1)$ (indices mod $4$), so $A=\{a_0,a_1,a_2,a_3\}$ and $B=A\cup\{v\}$. The critical values of $f$ are $a_1=f(c_1)$ and $v$. For an essential simple closed curve $\gamma\subset S^2\setminus B$ we write $T_\gamma$ for the Dehn twist about $\gamma$, an element of the pure mapping class group $\PMod(S^2,B)$, which acts on $\Teich(B)$ (see \S\ref{sec:teich}).

\begin{theorem}\label{thm:main}
\textbf{(a) Model case.} Let
\[f=f_{1/4}(z)=1-\frac{11}{12z}+\frac1{6z^2}.\]
Its $4$-cycle is $c_1=0\mapsto\infty\mapsto1\mapsto\tfrac14\mapsto0$, its second critical point is $c_2=4/11$, and $v=-25/96$. On $\Rh$ the points of $B$ are ordered as
\[v<c_1=0<f^{3}(c_1)=\tfrac14<f^{2}(c_1)=1<f(c_1)=\infty .\]
Let $\alpha$ be a simple arc from $v$ to $f(c_1)=\infty$ with the following crossing pattern with $\Rh$:
\begin{enumerate}[label=\textup{(\arabic*)}]
\item it leaves $v$ into the lower half-plane and crosses $\mathbb R$ upward at a point $x_1\in(1,\infty)$;
\item it runs in the upper half-plane and crosses $\mathbb R$ downward at a point $x_2\in(0,\tfrac14)$;
\item it runs in the lower half-plane and crosses $\mathbb R$ upward at a point $x_3\in(v,0)$;
\item it stays in the upper half-plane until $\infty$.
\end{enumerate}
Thus $\alpha$ meets $\Rh$ only at $v,x_1,x_2,x_3,\infty$, crossing transversally at $x_1,x_2,x_3$. A concrete choice is: the lower semicircle from $v$ to $2$; then the upper semicircle over $[\tfrac18,2]$; then the lower semicircle over $[-\tfrac18,\tfrac18]$; then the vertical ray from $-\tfrac18$ up to $\infty$.

Let $\gamma=\partial N(\alpha)$ be the boundary of a thin closed regular neighbourhood of $\alpha$. Then $\gamma$ is a simple closed curve in $S^2\setminus B$; it bounds a disc containing exactly $v$ and $f(c_1)$, and the complementary disc contains $c_1,f^2(c_1),f^3(c_1)$. The Dehn twist $T_\gamma$ itself (power $1$) satisfies
\[f^*\circ T_\gamma=f^*\qquad\text{and}\qquad i_*\circ T_\gamma=i_*,\]
so $T_\gamma(\D_f)=\D_f$. For every connected component $D_0$ of $\D_f$ and every $m\ne0$ we have $T_\gamma^m(D_0)\ne D_0$. Hence every orbit of $\langle T_\gamma\rangle$ on $\pi_0(\D_f)$ is infinite; in particular $\D_f$ has infinitely many components.

\smallskip
\textbf{(b) Every $f$.} Let $f$ be arbitrary as in the problem. By Proposition~\ref{prop:normal} there are a unique $t\in W$ and a unique M\"obius map $M$ with $M\circ f\circ M^{-1}=f_t$ and $M(c_1)=0$. Choose a path $t_s$ in $W$ from $\tfrac14$ to $t$ and let $h_1$ be the endpoint of an isotopy following the marked points (Lemma~\ref{lem:isotopy}). Then the curve
\[\gamma_f:=M^{-1}(h_1(\gamma))\]
has the same property for $f$: $T_{\gamma_f}$ preserves $\D_f$, and every orbit of $\langle T_{\gamma_f}\rangle$ on $\pi_0(\D_f)$ is infinite.

When $t\in\bigl(0,\tfrac{3-\sqrt5}{2}\bigr)$ this is explicit. The five points of $B$ lie on the circle $C=M^{-1}(\Rh)$ in the cyclic order $v,\,c_1,\,f^3(c_1),\,f^2(c_1),\,f(c_1)$. Call the \emph{upper} side of $C$ the image under $M^{-1}$ of the upper half-plane. Then $\gamma_f=\partial N(\alpha_f)$, where $\alpha_f$ is an arc from $v$ to $f(c_1)$ which
\begin{itemize}[nosep]
\item leaves $v$ into the lower side and crosses $C$ between $f^2(c_1)$ and $f(c_1)$;
\item crosses back between $c_1$ and $f^3(c_1)$;
\item crosses again between $v$ and $c_1$;
\item then runs on the upper side to $f(c_1)$.
\end{itemize}
\end{theorem}

\review{A sketch of the curve would be supremely helpful for exposition.}

\paragraph{What is known and what is proved here.} The setting of the problem, and the disconnectedness of $\D_f$ for quadratic maps with a $4$-periodic critical point, go back to E.~Hironaka and S.~Koch~\cite{HK}; for background on the pullback map in this generality see \cite{DH,BEKP}. The argument below is self-contained apart from the classical facts (F1)--(F5) listed in \S\ref{sec:teich}, which are quoted with references and not reproved. Everything else (the normal form, the fibre formula, the obstruction $(\star)$, the linking criterion, the face analysis at $t=\tfrac14$, the isotopy and lifting lemmas) is proved here.

\section{Overview of the proof}

After a M\"obius change of coordinates every $f$ is one of the maps $f_t(z)=1+b/z+c/z^2$ with cycle $0\to\infty\to1\to t\to0$, where $t$ ranges over a connected set $W$, the complement of nine points in $\mathbb C$. The space $\Teich(A)$ is the upper half-plane $\HH$, with the modular function $\lambda$ as the universal covering of the moduli space $\mathbb C\setminus\{0,1\}$, so the pair $\Psi=(i_*,f^*)$ maps $\Teich(B)$ to $\HH\times\HH$ and $\D_f=\Psi^{-1}(\text{diagonal})$.

The key algebraic fact is an explicit formula: if a point of $\Teich(B)$ has normalised coordinates $(0,\infty,1,t,v')$, then $w=\lambda(f^*\tau')$ satisfies $V_t(w)=v'$ for an explicit degree-$4$ rational function $V_t$. Since the fifth point $v'$ is never $\infty$, the value $w=1-t$ is never attained. Hence $\Psi(\Teich(B))$ avoids the graphs of all elements $g\in\Gamma(2)S$, where $S(\tau)=-1/\tau$; we call this the obstruction~$(\star)$. If $g$ has no fixed point in $\HH$, the complement $X_g$ of the graph of $g$ in $\HH\times\HH$ contains the diagonal and has $H_1(X_g)\cong\mathbb Z$, measured by winding about the graph. A mapping class $T$ with $\Psi\circ T=\Psi$ preserves $\D_f$. If, moreover, a path from $\tau$ to $T^m\tau$ has $\Psi$-image with nonzero class in $H_1(X_g)$, then $\tau$ and $T^m\tau$ cannot lie in the same component of $\D_f$, because a path inside $\D_f$ maps into the contractible diagonal. This is the linking criterion.

\review{Poor choice of notation and strategy to identify the Teichmuller space with the upper half-plane and use the modular function as the projection, and a congruence subgroup as the mapping class group throughout.}

\review{Confusion about identifying a point of the Teichmuller space with a quintuple of normalized coordinates, which is actually an element of the moduli space in the projection}

For $f_{1/4}$ we take $\gamma$ around the two critical values $v$ and $a_1$, along the arc $\alpha$. Then $T_\gamma$ is a point push of $v$, so $i_*\circ T_\gamma=i_*$, and its lift through $f$ is computed from the lift of $\alpha$ through $V_{1/4}$. The map $V_{1/4}$ is Galois with Klein four deck group, and its real locus consists of $\Rh$ and two circles; a face-by-face analysis shows that the lift $\beta$ of $\alpha$ starting at $w=\tfrac14$ ends at $\tfrac34$ after winding once clockwise around $0$. Lifting through $\lambda$ yields the element $g=T^{-2}S$, which is parabolic and so has no fixed point. From this we get $f^*\circ T_\gamma=f^*$, and the $\Psi$-image of the fibre path from $\tau_0$ to $T_\gamma^m\tau_0$ winds $\pm m$ times about the graph of $g$, which proves (a). Part (b) follows by transporting $\gamma$ along an isotopy that follows the marked points as $t$ moves in $W$, which conjugates all the structures involved.

\section{Normal form}\label{sec:normal}

\begin{proposition}\label{prop:normal}
Put $f_t(z)=1+\dfrac bz+\dfrac c{z^2}$, with $b=\dfrac{t^2+t-1}{1-t}$ and $c=\dfrac{t(1-2t)}{1-t}$. Let $E$ be the set of $9$ points
\[E=\{0,1,\tfrac12\}\cup\{t:t^2-3t+1=0\}\cup\{t:t^2+t-1=0\}\cup\{t:3t^2-3t+1=0\},\]
and put $W=\mathbb C\setminus E$. For $f$ as in the problem there is exactly one pair $(t,M)$ with $t\in W$, $M$ M\"obius, $MfM^{-1}=f_t$ and $M(c_1)=0$. Conversely, for every $t\in W$ the map $f_t$ is as in the problem, with cycle $0\to\infty\to1\to t\to0$. Finally, $W$ is connected.
\end{proposition}

\review{This is already worked out in (Firsova, Kahn, Selinger), which does not appear in the references. Also poor to define a space within a Proposition (in this case, $W$).}

\begin{proof}
\emph{Existence.} Let $M$ be the M\"obius map with $M(a_0,a_1,a_2)=(0,\infty,1)$, put $t=M(a_3)$ and $g=MfM^{-1}$. Then $0$ is a critical point of $g$ with $g(0)=\infty$, so $0$ is a double pole and the only pole; hence $g=P(z)/z^2$ with $\deg P\le2$ and $P(0)\ne0$. Since $g(\infty)=1$, we get $g=1+b/z+c/z^2$ with $c\ne0$. The conditions $g(1)=t$ and $g(t)=0$ read $1+b+c=t$ and $t^2+bt+c=0$. Substituting $c=t-1-b$ into the second gives $b(t-1)=1-t-t^2$. Hence $b$ and $c$ are as stated. (As checks, $b+c=t-1$ and $(1-t)(t^2+bt+c)=0$ identically.)

\emph{Uniqueness.} Any other such $M'$ maps $c_1=a_0$ to $0$, and then $a_1,a_2$ to $f_{t'}(0)=\infty$ and $f_{t'}(\infty)=1$; so $M'(a_0,a_1,a_2)=(0,\infty,1)$ and $M'=M$.

\emph{Which $t$ occur.} The value $t=\tfrac12$ gives $c=0$, and the degree drops. For $t\notin\{0,1,\tfrac12\}$ the map $f_t$ has degree $2$ and $0$ is a critical point of exact period $4$, since $0,\infty,1,t$ are distinct. Since $f_t'(z)=-(bz+2c)/z^3$, we have $c_2=-2c/b$, or $c_2=\infty=a_1$ if $b=0$, i.e.\ if $t^2+t-1=0$. Otherwise
\[v(t)=1-\frac{b^2}{4c}=-\frac{(t^2-3t+1)^2}{4t(1-t)(1-2t)},\]
using $(t^2+t-1)^2-(t^2-3t+1)^2=4t(1-t)(1-2t)$. Similarly, with $D=4t(1-t)(1-2t)$, one has $v-1=-(t^2+t-1)^2/D$ and $v-t=-(3t^2-3t+1)^2/D$; the latter because
\[(t^2-3t+1)^2+t\,D=9t^4-18t^3+15t^2-6t+1=(3t^2-3t+1)^2.\]
Hence $v=0\iff t^2-3t+1=0$, $v=1\iff b=0$, and $v=t\iff 3t^2-3t+1=0$. So $W$ is exactly the set of admissible parameters. It is $\mathbb C$ minus finitely many points, hence connected.
\end{proof}

At $t=\tfrac14$ one gets $b=-\tfrac{11}{12}$, $c=\tfrac16$, $c_2=\tfrac4{11}$ and $v=-\tfrac{25}{96}$, which gives the map $f_{1/4}$ of Theorem~\ref{thm:main}(a).

\paragraph{Factorisation.} Put $N_t(z)=z/(z-c_2(t))$ and $M_t(y)=v(t)+(1-v(t))/y$. Then
\begin{equation}\label{eq:factor}
f_t(z)=v+(1-v)\Bigl(1-\frac{c_2}{z}\Bigr)^2=M_t\bigl(N_t(z)^2\bigr).
\end{equation}
To check it, compare coefficients: $-2(1-v)c_2=b$ and $(1-v)c_2^2=c$, which follow from $1-v=b^2/(4c)$ and $c_2=-2c/b$. Both $M_t$ and $N_t$ depend continuously on $t\in W$. Also $N_t(0)=0$, $N_t(c_2)=\infty$, $M_t(0)=\infty$ and $M_t(\infty)=v$.

\section{Teichm\"uller preliminaries}\label{sec:teich}

\paragraph{Conventions.} A marking of $(S^2,X)$, with $|X|\ge3$, is an orientation-preserving homeomorphism $\varphi\colon S^2\to\Ch$. $\Teich(X)$ is the set of markings modulo the relation $\varphi_1\sim\varphi_2$, which holds iff there is a M\"obius map $M$ with $M\varphi_1=\varphi_2$ on $X$ and $M\varphi_1\simeq\varphi_2$ rel $X$. The pure mapping class group $\PMod(S^2,X)$ acts by $g\cdot[\varphi]=[\varphi g^{-1}]$; for a mapping class $T$ we write $F\circ T$ for the map $\tau\mapsto F(T\cdot\tau)$. The pullback $f^*[\varphi]=[\psi]$ is defined by requiring $\varphi\circ f\circ\psi^{-1}$ to be rational. In our situation $S^2=\Ch$, and $\tau_0=[\mathrm{id}]$ denotes the base point.

\paragraph{Standard facts quoted, not reproved.}
\begin{enumerate}[label=\textbf{(F\arabic*)},leftmargin=*]
\item \label{F1} (Modular function; classical, see e.g.\ \cite{Ahlfors,Chandra}.) Let $\lambda=(\theta_2/\theta_3)^4\colon\HH\to\mathbb C\setminus\{0,1\}$. It is a holomorphic universal covering whose fibres are the orbits of $\Gamma(2)$. Moreover $\lambda(\tau+2)=\lambda(\tau)$ and $\lambda(-1/\tau)=1-\lambda(\tau)$; $\lambda(\tau)=16e^{i\pi\tau}+O(e^{2\pi i\tau})$ as $\Img\tau\to\infty$; and $y\mapsto\lambda(iy)$ is a decreasing bijection $(0,\infty)\to(0,1)$.
\item \label{F2} (Teichm\"uller; see e.g.\ \cite{Hubbard}.) $\Teich(X)$ is a complex manifold homeomorphic to $\mathbb R^{2(|X|-3)}$; in particular it is contractible. $\PMod(S^2,X)$ acts properly discontinuously by biholomorphisms, and the quotient is the moduli space $\M_{0,X}$.
\item \label{F3} (Bers fibration and Birman exact sequence \cite{Bers,Birman,FM}.) The forgetful map $i_*\colon\Teich(B)\to\Teich(A)$ is surjective. Fix $p\in\Teich(A)$ and normalise $a_0,a_1,a_2,a_3$ to $0,\infty,1,t$; write $A_p=\{0,\infty,1,t\}$. Recording the position $v'$ of the fifth point identifies the fibre $i_*^{-1}(p)$ with the universal cover of $\Ch\setminus A_p$. The point-push subgroup $\ker\bigl(\PMod(S^2,B)\to\PMod(S^2,A)\bigr)\cong\pi_1(S^2\setminus A,v)$, $[\ell]\mapsto\Push(\ell)$, acts on each fibre as the deck group: for the fibre through $\tau_0$, $\Push(\ell)^{\pm1}\cdot\tau_0$ is the endpoint of the lift through $v'$ of $\ell$ starting at $\tau_0$, the sign depending only on conventions.
\item \label{F4} (Push formula, see \cite{FM}.) The push of $v$ along a simple loop $\delta$ is $T_{\delta_L}T_{\delta_R}^{-1}$, where $\delta_L,\delta_R$ are the two boundary curves of a regular neighbourhood of $\delta$.
\item \label{F5} (Douady--Hubbard \cite{DH}.) $f^*$ is well defined and continuous; indeed it is holomorphic. (In \cite{DH} this is proved for Thurston's pullback map; the proof applies verbatim to $f^*\colon\Teich(B)\to\Teich(A)$.)
\end{enumerate}
The push formula is used only up to sign, and every argument below is symmetric under $m\mapsto-m$.

\begin{lemma}[free action]\label{lem:free}
$\PMod(S^2,X)$ acts freely on $\Teich(X)$.
\end{lemma}

\review{This is classical.}

\begin{proof}
Suppose $[\varphi g^{-1}]=[\varphi]$. Then $M\varphi=\varphi g^{-1}$ on $X$ for some M\"obius $M$. Since $g$ fixes $X$ pointwise, $M$ fixes the at least three points of $\varphi(X)$, so $M=\mathrm{id}$. Then $\varphi g^{-1}\simeq\varphi$ rel $X$, so $g=1$ in $\PMod(S^2,X)$.
\end{proof}

By \ref{F2} and Lemma~\ref{lem:free}, the modulus map $\mathrm{mod}\colon\Teich(A)\to\mathbb C\setminus\{0,1\}$ is a covering, and since $\Teich(A)$ is simply connected it is a universal covering. It sends $[\varphi]$ to the $t$ for which some M\"obius map takes $\varphi(a_0,a_1,a_2,a_3)$ to $(0,\infty,1,t)$. By \ref{F1} there is a biholomorphism $\Theta\colon\Teich(A)\to\HH$ with $\lambda\circ\Theta=\mathrm{mod}$, and we use it to identify $\Teich(A)$ with $\HH$. For $f=f_{1/4}$ the base point $\tau_0=[\mathrm{id}]\in\Teich(B)$ lies in $\D_f$, because $f^*[\mathrm{id}]=[\mathrm{id}]$, $f$ being rational. Composing $\Theta$ with an element of $\Gamma(2)$ we normalise it so that $i_*\tau_0=p_0:=iy_0$, where $\lambda(iy_0)=\tfrac14$ (numerically $y_0\approx1.2793$).

Put $\Psi=(i_*,f^*)\colon\Teich(B)\to\HH\times\HH$. Then $\D_f=\Psi^{-1}(\Delta)$, where $\Delta$ is the diagonal.

\begin{lemma}[equivariance for a lift fixing $A$]\label{lem:equiv}
Let $\varphi$ be a homeomorphism fixing $B$ pointwise, and $\tilde\varphi$ a homeomorphism fixing $A$ pointwise with $f\tilde\varphi=\varphi f$. Then $f^*(\varphi\cdot\tau)=\tilde\varphi\cdot f^*(\tau)$.
\end{lemma}

\review{Also classical.}

\begin{proof}
Let $\tau=[\psi_0]$ and suppose $\psi_0 f=F\psi$ with $F$ rational. Then $\psi_0\varphi^{-1}f=\psi_0 f\tilde\varphi^{-1}=F(\psi\tilde\varphi^{-1})$, so $f^*[\psi_0\varphi^{-1}]=[\psi\tilde\varphi^{-1}]$.
\end{proof}

A lift of an element of $\PMod(S^2,B)$ need not fix $A$ pointwise: it may exchange $a_i$ with the other preimage of $a_{i+1}$. So Lemma~\ref{lem:equiv} has to be applied to an explicitly constructed lift, as in the next lemma.

\begin{lemma}[lift of $T_\gamma$]\label{lem:lift}
Let $\alpha$ be a simple arc from $v$ to $a_1$ with interior in $S^2\setminus B$. Let $\gamma=\partial N(\alpha)$, and realise $T_\gamma$ in a collar $R\subset S^2\setminus N(\alpha)$ disjoint from $B$. Then:
\begin{enumerate}[label=\textup{(\roman*)},nosep]
\item $f^{-1}(R)=R_1\sqcup R_2$, with each $R_j$ mapped homeomorphically onto $R$;
\item the map $\tilde T$ defined as $(f|_{R_j})^{-1}\,T_\gamma\, f$ on $R_j$ and the identity elsewhere is a homeomorphism with $f\tilde T=T_\gamma f$;
\item $\tilde T$ fixes $f^{-1}(B)\supset A$ pointwise.
\end{enumerate}
Hence $f^*\circ T_\gamma=[\tilde T]\cdot f^*$ with $[\tilde T]\in\PMod(S^2,A)$.
\end{lemma}

\begin{proof}
$R$ contains no critical value, so $f^{-1}(R)\to R$ is a $2$-sheeted covering. Its monodromy around the core of $R$ is the product of the monodromies around the two critical values $v,a_1$ inside $\gamma$. Each of these is the transposition of the two sheets, so the product is trivial. Hence the covering splits into two annuli, each mapped homeomorphically. Since $T_\gamma=\mathrm{id}$ on $\partial R$, $\tilde T$ is continuous and bijective, hence a homeomorphism of $S^2$, and $f\tilde T=T_\gamma f$ holds on each $R_j$ and trivially elsewhere. The set $f^{-1}(B)$ is disjoint from $f^{-1}(R)$, so $\tilde T$ fixes it; it contains $A$ because $f(A)=A\subset B$. Lemma~\ref{lem:equiv} gives the last statement.
\end{proof}

\section{The fibre formula and the obstruction \texorpdfstring{$(\star)$}{(*)}}

\begin{lemma}\label{lem:fibreformula}
Let $\tau'\in\Teich(B)$ be represented, after normalising, by
\[(\Ch;a_0,a_1,a_2,a_3,v)=(0,\infty,1,t,v').\] Then $w:=\lambda(f^*\tau')$ satisfies
\[V_t(w)=v',\qquad\text{where}\quad V_t(w)=-\frac{\bigl((w-1)^2-t\bigr)^2}{4w(w-1)(w+t-1)} .\]
\end{lemma}

\review{The aforementioned representation of an element of the Teichmuller space is actually an element of the moduli space.}

\begin{proof}
The pulled-back structure carries a degree-$2$ rational map $F$ whose critical values are the images $\infty$ and $v'$ of the critical values $a_1$ and $v$. Up to a M\"obius change in the source, $F(z)=v'+z^2$. Since $a_0=c_1$ is the critical point over $a_1$, its marked image is $a_0'=\infty$. The other points satisfy $F(a_i')=$ (image of $a_{i+1}$), so $a_1'=s_2$, $a_2'=s_3$, $a_3'=s_0$ with
\[s_2^2=1-v',\qquad s_3^2=t-v',\qquad s_0^2=-v';\]
the marking determines the signs. The map $z\mapsto(s_3-s_2)/(z-s_2)$ sends $(\infty,s_2,s_3,s_0)$ to $(0,\infty,1,w)$, so $w=(s_3-s_2)/(s_0-s_2)$.

Put $u=s_2-s_0$ and $U=u^2$. From $s_2^2-s_0^2=1$ we get $s_2+s_0=u^{-1}$, so $s_2=(u+u^{-1})/2$ and $s_0=(u^{-1}-u)/2$. Also $s_3=s_2-wu$. The relation $s_3^2-s_0^2=t$ becomes $1-w(U+1)+w^2U=t$, so
\[U=\frac{w+t-1}{w(w-1)},\qquad 1-U=\frac{(w-1)^2-t}{w(w-1)}.\]
Then $v'=-s_0^2=-(1-U)^2/(4U)$, which is the stated formula.
\end{proof}

Note that $V_t(t)=-\dfrac{(t^2-3t+1)^2}{4t(t-1)(2t-1)}=v(t)$.

\begin{corollary}[the obstruction $(\star)$]\label{cor:star}
$V_t^{-1}(\infty)=\{0,1,\infty,1-t\}$, and at $w=1-t$ the numerator of $V_t$ equals $(t^2-t)^2\ne0$. Since $v'\ne a_1=\infty$, we get $w\ne1-t$. Thus
\[\lambda(f^*\tau')\ne1-\lambda(i_*\tau')\qquad\text{for all }\tau'\in\Teich(B). \tag{$\star$}\]
By \ref{F1}, $\lambda(q)=1-\lambda(p)=\lambda(-1/p)$ holds iff $q=g(p)$ for some $g\in\Gamma(2)S$, where $S(\tau)=-1/\tau$. So $\Psi(\Teich(B))$ avoids every graph $\Gamma_g=\{(p,g(p)):p\in\HH\}$ with $g\in\Gamma(2)S$.
\end{corollary}

\begin{proof}
The denominator vanishes at $0,1,1-t$, which are distinct since $t\ne0,1$, and the numerator is $(1-t)^2$, $t^2$, $(t^2-t)^2$ there, all nonzero. Since the numerator has degree $4$ and the denominator degree $3$, also $V_t(\infty)=\infty$. The rest was explained in the statement, using $\lambda(i_*\tau')=t$.
\end{proof}

\begin{lemma}[branching of $V_t$]\label{lem:branch}
For $t\in W$, $V_t$ has degree $4$ and is unbranched over $\Ch\setminus\{0,1,t\}$.
\end{lemma}

\begin{proof}
By the proof of Corollary~\ref{cor:star}, numerator and denominator are coprime, so $\deg V_t=4$. Fix $v'\notin\{0,1,t,\infty\}$. If $V_t(w)=v'$, then $w\notin\{0,1,1-t,\infty\}$; reversing the computation in the proof of Lemma~\ref{lem:fibreformula} (choose $u$ with $u^2=U$ and define $s_2,s_0,s_3$ by the formulas there) shows that $w=(s_3-s_2)/(s_0-s_2)$ for some choice of square roots $s_2,s_3,s_0$ of $1-v',t-v',-v'$. So the solutions $w$ come from the sign choices of $(s_2,s_3,s_0)$ modulo a global sign, and there are at most $4$ of them.

They are distinct. If two sign classes gave the same $w$, the M\"obius map carrying one quadruple $(\infty,s_2,s_3,s_0)$ to the other would be an affine map $L$ with $L(s_j)=\pm s_j$, with a sign pattern that is not constant. Then $L(z)^2-z^2$, of degree at most $2$, would vanish at the three distinct points $s_2,s_3,s_0$, forcing $L=\pm\mathrm{id}$, a contradiction. Over $\infty$ there are $4$ simple preimages $0,1,1-t,\infty$. Hence every point outside $\{0,1,t\}$ has $4$ distinct preimages.
\end{proof}

\begin{lemma}[fibre lemma]\label{lem:fibre}
Fix $p\in\Teich(A)$ with $\lambda(p)=t$, and let $\tilde U=i_*^{-1}(p)$. By \ref{F3}, $\tilde U$ is the universal cover of $\Ch\setminus A_p$ via $\tau'\mapsto v'(\tau')$. Then $w=\lambda\circ f^*$ is the lift of $v'$ through the covering $V_t\colon\Ch\setminus V_t^{-1}(A_p)\to\Ch\setminus A_p$, and $f^*|_{\tilde U}$ is the lift of $w$ through $\lambda$. Both lifts are determined by their value at one point.
\end{lemma}

\begin{proof}
By Lemma~\ref{lem:branch}, $V_t$ restricted as stated is a covering, since all its critical values lie in $A_p$. By \ref{F5} and Lemma~\ref{lem:fibreformula}, $w=\lambda\circ f^*$ is continuous on $\tilde U$ and satisfies $V_t\circ w=v'$; as $v'\notin A_p$, $w$ takes values in $\Ch\setminus V_t^{-1}(A_p)$. So $w$ is a lift of $v'$, and $f^*|_{\tilde U}$ is a lift of $w$ through $\lambda$ by definition. Since $\tilde U$ is connected, lifts are determined by their value at one point.
\end{proof}

\section{\texorpdfstring{$\D_f$}{D\_f} is a smooth curve}

\begin{lemma}\label{lem:smooth}
$\D_f$ is a one-dimensional complex submanifold of $\Teich(B)$. Hence its connected components are path components.
\end{lemma}

\review{Known: proved by Epstein, which does not appear in the references.}

\begin{proof}
The map $\Teich(B)\to\M_{0,5}$, $\tau'\mapsto(t',v')$, is a covering by \ref{F2} and Lemma~\ref{lem:free}, so $(t',v')$ are local coordinates. Let $\tau_1\in\D_f$ have coordinates $(t_1,v_1)$. Then $w(\tau_1)=\lambda(f^*\tau_1)=\lambda(i_*\tau_1)=t_1$, so $v_1=V_{t_1}(t_1)=v(t_1)$. Since $v_1\notin\{0,1,t_1,\infty\}$, Lemma~\ref{lem:branch} makes $V_{t_1}$ a local biholomorphism at $w=t_1$. By the implicit function theorem and continuity of $w$, near $\tau_1$ the function $w(t',v')$ is the holomorphic branch of $V_{t'}^{-1}(v')$ through $t_1$. Moreover $\lambda$ is injective near $i_*\tau_1=f^*\tau_1$. So near $\tau_1$ the condition $f^*\tau'=i_*\tau'$ is equivalent to $w(t',v')=t'$, i.e.\ to $v'=V_{t'}(t')=v(t')$, which defines a holomorphic graph.
\end{proof}

\section{The linking criterion}

\begin{proposition}\label{prop:linking}
Let $T\in\PMod(S^2,B)$ satisfy $\Psi\circ T=\Psi$, so that $T^{\pm1}(\D_f)=\D_f$. Let $g\in\Gamma(2)S$ have no fixed point in $\HH$. Let $\tau\in\D_f$ with $p_0=i_*\tau$, and let $D_0$ be its component. Suppose that for each $m\ne0$ there is a path $\sigma'_m$ in the fibre $i_*^{-1}(p_0)$ from $\tau$ to $T^m\tau$ such that the loop $f^*\circ\sigma'_m$ has nonzero winding number about $g(p_0)$. Then $T^m D_0\ne D_0$ for every $m\ne0$.
\end{proposition}

\begin{proof}
Let $X_g=\HH\times\HH\setminus\Gamma_g$. The map $(p,q)\mapsto\bigl(p,(q-\Rea g(p))/\Img g(p)\bigr)$ identifies $X_g$ with $\HH\times(\HH\setminus\{i\})$, and is an orientation-preserving affine map on each fibre $\{p\}\times\HH$. So $H_1(X_g)\cong\mathbb Z$, and the class of a loop in $\{p\}\times\HH$ is its winding number about $g(p)$. By $(\star)$, $\Psi(\Teich(B))\subset X_g$, and since $g$ has no fixed point, $\Delta\subset X_g$.

Suppose $T^mD_0=D_0$. By Lemma~\ref{lem:smooth} there is a path $\sigma\subset D_0$ from $\tau$ to $T^m\tau$. Since $\Teich(B)$ is contractible \ref{F2}, $\sigma\simeq\sigma'_m$ rel endpoints. Applying $\Psi$, the loops $\Psi\sigma$ and $\Psi\sigma'_m$ (they are loops because $\Psi(T^m\tau)=\Psi(\tau)$) are homotopic in $X_g$. But $\Psi\sigma\subset\Delta$, which is contractible, so $\Psi\sigma$ is null-homologous in $X_g$, whereas $\Psi\sigma'_m=(p_0,f^*\sigma'_m)$ has nonzero class. This is a contradiction.
\end{proof}

\section{The computation at \texorpdfstring{$t_0=\tfrac14$}{t0=1/4}}\label{sec:comp}

From now on $f=f_{1/4}$ and $V=V_{1/4}$. Since $(w-1)^2-\tfrac14=(w-\tfrac12)(w-\tfrac32)$,
\[V(w)=-\frac{(w-\frac12)^2(w-\frac32)^2}{4w(w-1)(w-\frac34)} .\]

\paragraph{Deck group.} Put $g_0(w)=\dfrac{w-3/4}{w-1}$, $g_2(w)=\dfrac3{4w}$ and $g_3=g_0g_2=g_2g_0=\dfrac{3(w-1)}{4(w-3/4)}$; $g_0$ and $g_2$ are commuting involutions. Direct substitution gives $V\circ g_0=V\circ g_2=V$:
\begin{itemize}[nosep]
\item under $g_2$: $w-\tfrac12\mapsto-(w-\tfrac32)/(2w)$, $w-\tfrac32\mapsto-3(w-\tfrac12)/(2w)$, and $4w(w-1)(w-\tfrac34)\mapsto 9(w-\tfrac34)(w-1)/(4w^3)$;
\item under $g_0$: $w-\tfrac12\mapsto(w-\tfrac12)/(2(w-1))$, $w-\tfrac32\mapsto-(w-\tfrac32)/(2(w-1))$, and the denominator becomes $w(w-\tfrac34)/(4(w-1)^3)$.
\end{itemize}
So $V$ is Galois with Klein four deck group $\{\mathrm{id},g_0,g_2,g_3\}$. Its $6$ critical points are the fixed points of $g_0,g_2,g_3$:
$\tfrac12,\tfrac32$, which lie over $0$; $\pm\tfrac{\sqrt3}2$, which lie over $1$; and $P,\bar P=\tfrac34\pm\tfrac{\sqrt3}4 i$, which lie over $\tfrac14$. The critical values are therefore $0,1,\tfrac14\in A$.

\paragraph{Real locus.} $V$ has real coefficients, so $V(w)\in\Rh$ iff $V(\bar w)=V(w)$ iff $\bar w=h(w)$ for some deck map $h$:
$h=\mathrm{id}$ gives $\Rh$; $h=g_0$ gives the circle $C_1\colon|w-1|=\tfrac12$; $h=g_2$ gives $C_2\colon|w|=\tfrac{\sqrt3}2$; and $h=g_3$ gives $|w-\tfrac34|^2=-\tfrac3{16}$, which is empty. Moreover $C_1\cap C_2=\{P,\bar P\}$.

\paragraph{Faces.} All critical values of $V$ are real. So $V$ restricted to $\Ch\setminus V^{-1}(\Rh)$ is a covering of the two open half-planes, and each complementary component (face) maps homeomorphically onto one half-plane. Adjacent faces map to opposite half-planes, because $V$ is a local homeomorphism along each edge. We need four faces, all inside $C_2$ and hence avoiding $w=1$:
\begin{gather*}
L_2=\{\Img w<0,\ |w|<\tfrac{\sqrt3}2,\ |w-1|>\tfrac12\},\qquad U_2=\overline{L_2}\ \text{(complex conjugate)},\\
U_{12}=\{\Img w>0,\ |w|<\tfrac{\sqrt3}2,\ |w-1|<\tfrac12\},\qquad L_{12}=\overline{U_{12}}.
\end{gather*}

\paragraph{Boundary values.} The vertices of the graph $V^{-1}(\Rh)$ are the points of $V^{-1}(\{0,1,\tfrac14,\infty\})=\{\tfrac12,\tfrac32,\pm\tfrac{\sqrt3}2,P,\bar P,0,1,\tfrac34,\infty\}$. Each edge contains no vertex except at its ends, so it maps homeomorphically onto the arc of $\Rh$ between the images of its ends that avoids the other marked values.
\begin{itemize}[nosep]
\item On $(-\tfrac{\sqrt3}2,0)$, $V$ increases from $1$ to $+\infty$: here $V\approx-3/(16w)$ near $0$, and $V(-\tfrac{\sqrt3}2)=1$.
\item On $(0,\tfrac12)$, $V$ increases from $-\infty$ to $0$, with $V(\tfrac14)=v$.
\item On $(\tfrac12,\tfrac34)$, $V$ decreases from $0$ to $-\infty$, with $V(\tfrac23)=v$, since $\tfrac23=g_0(\tfrac14)$.
\item The $C_1$-arcs from $\tfrac12$ to $P$ and from $\tfrac12$ to $\bar P$ map onto $(0,\tfrac14)$: the alternative arc of $\Rh$ from $0$ to $\tfrac14$ would pass through $\infty$ and $1$, whose preimages $0,1,\tfrac34,\pm\tfrac{\sqrt3}2$ are not on these arcs.
\item The $C_2$-arcs from $P$ to $-\tfrac{\sqrt3}2$ (outside $C_1$) and from $P$ to $\tfrac{\sqrt3}2$ (inside $C_1$), and their conjugates, map onto $(\tfrac14,1)$; the segment $(\tfrac34,\tfrac{\sqrt3}2)$ maps onto $(1,\infty)$.
\item $V'(\tfrac14)>0$ and $V$ is real on $\mathbb R$, so $U_2$, which lies just above $\tfrac14$, maps to the upper half-plane.
\end{itemize}
Since $U_2$ is adjacent to $L_2$ across $(-\tfrac{\sqrt3}2,0)$ and to $U_{12}$ across the $C_1$-arc from $\tfrac12$ to $P$, and $U_{12}$ is adjacent to $L_{12}$ across $(\tfrac12,\tfrac34)$, it follows that $U_2$ and $L_{12}$ map to the upper half-plane, and $L_2$ and $U_{12}$ to the lower.

\paragraph{Lift of $\alpha$.} Let $\beta$ be the lift of $\alpha$ through $V$ starting at $w=\tfrac14$, which lies over $v$. The four sub-arcs of $\alpha$ between consecutive points of $\alpha\cap\Rh$ each lie in one open half-plane, and on each face the inverse of $V$ is a homeomorphism onto the closure of that face; at the crossing points $V$ is a local biholomorphism, real on $\mathbb R$.
\begin{enumerate}[label=\textup{\arabic*.},nosep]
\item $\alpha$ leaves $v$ downward. Since $V'(\tfrac14)>0$, $\beta$ enters $L_2$. On $\partial L_2$ the preimage of $(1,\infty)$ is $(-\tfrac{\sqrt3}2,0)$, where $V$ is increasing, so $\beta$ crosses $\mathbb R$ there, upward, into $U_2$.
\item On $\partial U_2$ the preimage of $(0,\tfrac14)$ is the $C_1$-arc from $\tfrac12$ to $P$. $\beta$ crosses it into $U_{12}$.
\item On $\partial U_{12}$ the preimage of $(v,0)$ is $(\tfrac12,\tfrac23)$, where $V$ is decreasing. $\beta$ crosses it downward into $L_{12}$.
\item The only preimage of $\infty$ in $\overline{L_{12}}$ is $\tfrac34$ (the poles $0$ and $1$ lie outside $C_1$ and $C_2$ respectively). So $\beta$ ends at $w=\tfrac34=1-t_0$.
\end{enumerate}
In particular only the crossing pattern of $\alpha$ was used, not the concrete realisation. The arc $\beta$ lies in the disc $|w|<\tfrac{\sqrt3}2$, which avoids $1$, and avoids $0$. It meets $\mathbb R$ only at $\tfrac14$, at one point of the negative real axis (crossed upward), at one point of $(\tfrac12,\tfrac23)$ (crossed downward), and at $\tfrac34$, which it reaches from below. Following the argument of $w$ along $\beta$ followed by the segment $[\tfrac34,\tfrac14]$, it decreases by $2\pi$. So this closed loop has winding number $-1$ about $0$, and since the homotopy class of a loop in a punctured disc is determined by its winding number, in $\mathbb C\setminus\{0,1\}$
\begin{equation}\label{eq:beta}
\beta\simeq(\text{clockwise loop about }0\text{ based at }\tfrac14)\cdot[\tfrac14,\tfrac34]\qquad\text{rel endpoints}.
\end{equation}

\paragraph{Lift through $\lambda$.} The clockwise loop about $0$ is homotopic to: go from $\tfrac14$ to $\varepsilon$ along $\mathbb R$, around the small clockwise circle $|w|=\varepsilon$, and back. By \ref{F1} its $\lambda$-lift from $p_0=iy_0$ goes up the imaginary axis, then decreases $\Rea\tau$ by $2$ (as $\lambda\approx16e^{i\pi\tau}$ for large $\Img\tau$), and comes back down; it ends at $p_0-2$. By $\lambda(-1/\tau)=1-\lambda(\tau)$ and $\lambda(\tau-2)=\lambda(\tau)$, the segment $[\tfrac14,\tfrac34]$ then lifts to the vertical segment from $p_0-2$ to $-2+i/y_0$. So the endpoint of the $\lambda$-lift of $\beta$ from $p_0$ is
\[q^*=g(p_0),\qquad g=T^{-2}S=\begin{pmatrix}-2&-1\\1&0\end{pmatrix}\in\Gamma(2)S,\qquad \operatorname{tr}g=-2,\]
where $T$ here denotes $\tau\mapsto\tau+1$. So $g$ is parabolic and has no fixed point in $\HH$.

\begin{remark}[orientation conventions]
Had the winding of $\beta$ about $0$ been $+1$ instead of $-1$ (e.g.\ under the opposite orientation convention), the same computation would give $g=T^{2}S=\begin{pmatrix}2&-1\\1&0\end{pmatrix}$, of trace $2$, again parabolic. More generally $T^{2k}S$ has trace $2k$ and is fixed-point-free on $\HH$ for every $k\ne0$; only the fact that the winding is nonzero matters.
\end{remark}

\section{Proof of Theorem~\ref{thm:main}(a)}\label{sec:proofa}

\paragraph{$\alpha$ is simple and $\gamma$ is as claimed.} In the concrete realisation, the lower semicircle from $v$ to $2$ has centre $\approx0.870$ and radius $\approx1.130$; the small lower semicircle over $[-\tfrac18,\tfrac18]$ lies within distance $0.995$ of that centre, so strictly inside it. The upper semicircle over $[\tfrac18,2]$ has real parts in $[\tfrac18,2]$, so it misses the vertical ray at $\Rea z=-\tfrac18$. The arc meets $\Rh$ only at $v,2,\tfrac18,-\tfrac18,\infty$, with $2\in(1,\infty)$, $\tfrac18\in(0,\tfrac14)$ and $-\tfrac18\in(v,0)$, crossing transversally in the required directions. So $\alpha$ is simple, has the crossing pattern, and meets $B$ only at its ends. For any simple arc $\alpha$ with the crossing pattern, $\gamma=\partial N(\alpha)$ bounds a disc containing exactly $v,a_1$, because $\alpha$ is simple with interior in $S^2\setminus B$; and the lifting argument of \S\ref{sec:comp} used only the crossing pattern. Everything below therefore applies to any such arc.

\paragraph{$T_\gamma$ is a point push.} Let $\ell$ be the lasso: follow $\alpha$ to near $a_1$, go once around $a_1$, and return. A regular neighbourhood of $\ell$ is an annulus. One of its boundary curves bounds a disc containing only $a_1$ (so the twist about it is trivial), and the other is $\gamma$. So by \ref{F4}, $\Push(\ell)=T_\gamma^{\pm1}$. In particular $T_\gamma$ maps to $1$ in $\PMod(S^2,A)$, so $i_*\circ T_\gamma=i_*$.

\paragraph{$T_\gamma$ preserves $\D_f$.} Work in the fibre $i_*^{-1}(p_0)$ through $\tau_0$; here $A_{p_0}=\{0,\infty,1,\tfrac14\}$, $v'(\tau_0)=v$, $w(\tau_0)=\lambda(f^*\tau_0)=\tfrac14$ and $f^*\tau_0=p_0$, because $\tau_0\in\D_f$.
\begin{itemize}[nosep]
\item By \ref{F3}, there is a path $\sigma'$ from $\tau_0$ to $T_\gamma^{\pm1}\tau_0$ which projects to $\ell$ under $v'$.
\item Truncate $\alpha$ to $\alpha'$, stopping just before $a_1$. By the fibre lemma~\ref{lem:fibre}, $w\circ\sigma'$ is the $V$-lift of $\ell$ starting at $\tfrac14$, which is $\beta'\cdot c\cdot\beta'^{-1}$. Here $\beta'$ is the corresponding truncation of $\beta$, and $c$ is a small loop about $\tfrac34$; this uses that $\tfrac34$ is a simple preimage of $\infty$.
\item Since $\tfrac34\notin\{0,1\}$, $c$ is null-homotopic in $\mathbb C\setminus\{0,1\}$. So its $\lambda$-lift is a small closed loop $\tilde c$ about $q^*$ (recall $\lambda(q^*)=\tfrac34$ and $\lambda$ is a local biholomorphism), and $f^*\sigma'=\tilde\beta'\cdot\tilde c\cdot\tilde\beta'^{-1}$ is closed, where $\tilde\beta'$ is the $\lambda$-lift of $\beta'$ from $p_0$.
\item Hence $f^*(T_\gamma^{\pm1}\tau_0)=f^*(\tau_0)$. By Lemma~\ref{lem:lift}, $[\tilde T]^{\pm1}$ fixes the point $f^*\tau_0$, so $[\tilde T]=1$ by Lemma~\ref{lem:free}.
\item Therefore $f^*\circ T_\gamma=f^*$ on all of $\Teich(B)$. Together with $i_*\circ T_\gamma=i_*$ this gives $\Psi\circ T_\gamma=\Psi$, and so $T_\gamma^{\pm1}(\D_f)=\D_f$.
\end{itemize}

\paragraph{Infinite orbits.} Fix $m\ne0$, and let $\sigma'_m$ be the path in the fibre from $\tau_0$ to $T_\gamma^m\tau_0$ lying over the loop $\ell^{\pm m}$ (\ref{F3}). As above, $f^*\sigma'_m=\tilde\beta'\,\tilde c^{\,\pm m}\,\tilde\beta'^{-1}$. This loop lies in $\HH\setminus\{g(p_0)\}$ by $(\star)$, and its winding number about $q^*=g(p_0)$ is $\pm m\ne0$, because $\tilde c$ winds once about $q^*$ (it is the image of a small circle about $\tfrac34$ under a local inverse of $\lambda$). So the class of $\Psi\circ\sigma'_m$ in $H_1(X_g)\cong\mathbb Z$ is $\pm m$, and Proposition~\ref{prop:linking} (with $T=T_\gamma$, $\tau=\tau_0$) gives $T_\gamma^mD_0\ne D_0$ for the component $D_0$ of $\tau_0$.

Now let $\tau_1\in\D_f$ be arbitrary and $D_1$ its component. Choose a path $\rho$ from $\tau_0$ to $\tau_1$. Then $\rho^{-1}\cdot\sigma'_m\cdot T_\gamma^m\rho$ is a path from $\tau_1$ to $T_\gamma^m\tau_1$. Since $\Psi\circ T_\gamma^m=\Psi$, its $\Psi$-image is the conjugate loop $(\Psi\rho)^{-1}(\Psi\sigma'_m)(\Psi\rho)$, which lies in $X_g$ and has class $\pm m$. Suppose $T_\gamma^mD_1=D_1$, and let $\sigma\subset D_1$ be a path from $\tau_1$ to $T_\gamma^m\tau_1$ (Lemma~\ref{lem:smooth}). Contractibility of $\Teich(B)$ makes the two paths homotopic rel endpoints, so their $\Psi$-images are homotopic in $X_g$. But $\Psi\sigma\subset\Delta$, which is contractible, so its class is $0$, a contradiction. This is the argument of Proposition~\ref{prop:linking} with the base point moved to $\tau_1$. Hence $T_\gamma^mD_1\ne D_1$ for every component $D_1$ and every $m\ne0$: every orbit of $\langle T_\gamma\rangle$ on $\pi_0(\D_f)$ is infinite. This proves Theorem~\ref{thm:main}(a). \qed

\begin{remark}[a direct topological proof that the lift is trivial]
The identity $f^*\circ T_\gamma=f^*$ can also be seen without the fibre computation. Put $s=1/z-\tfrac{11}4$. Then $f=v+s^2/6$, since $v+\tfrac16(1/z-\tfrac{11}4)^2=v+\tfrac{121}{96}-\tfrac{11}{12z}+\tfrac1{6z^2}$ and $v+\tfrac{121}{96}=1$. The critical points are $s=0$ ($z=c_2$, over $v$) and $s=\infty$ ($z=0=a_0$, over $a_1=\infty$). In this coordinate $s(a_0)=\infty$, $s(a_1)=-\tfrac{11}4$, $s(a_2)=-\tfrac74$, $s(a_3)=\tfrac54$, and the positive numbers $\tfrac54<\tfrac74<\tfrac{11}4$ satisfy $v+s^2/6=0,\tfrac14,1$ respectively.

Put $Y=f-v=s^2/6$. The arc $\alpha$ meets $\Rh$, apart from its endpoints, only at $x_1,x_2,x_3>v$, i.e.\ at positive $Y$. So the interior of $\alpha$ lies in $\mathbb C\setminus(-\infty,0]$ in the $Y$-coordinate, where $s=\sqrt{6Y}$ has a branch with $\Rea s>0$. Hence $f^{-1}(\alpha)=\alpha_+\cup(-\alpha_+)$ with $\alpha_+\subset\{\Rea s>0\}$ a simple arc from $s=0$ to $s=\infty$, and $J=f^{-1}(\alpha)$ is a Jordan curve through $0$ and $\infty$. The involution $s\mapsto-s$ preserves $J$ and exchanges its two complementary discs $\Omega,\Omega'$: near its fixed point $0\in J$ it is a rotation by $\pi$ which exchanges the two local sides of $J$.

$J$ meets the positive real $s$-axis only in $\alpha_+$, at the three transversal crossings $s_i=\sqrt{6(x_i-v)}$, and $x_3<0<x_2<\tfrac14<1<x_1$ gives $s_3<\tfrac54<s_2<\tfrac74<\tfrac{11}4<s_1$. So $\tfrac54$ and $\tfrac74$ lie on opposite sides of $J$, and $\tfrac74,\tfrac{11}4$ on the same side. Say $\tfrac54\in\Omega$; then $\tfrac74,\tfrac{11}4\in\Omega'$, so $-\tfrac74,-\tfrac{11}4\in\Omega$. Hence $a_1,a_2,a_3\in\Omega$ and $a_0\in J$. Now $f^{-1}(N(\alpha))$ is a regular neighbourhood of $J$, and the collars $R_1,R_2$ of Lemma~\ref{lem:lift} are parallel to its two boundary curves $\gamma_1\subset\Omega$, $\gamma_2\subset\Omega'$. The curve $\gamma_1$ separates $\{a_1,a_2,a_3\}$ from $\{a_0\}$, so it is peripheral in $(S^2,A)$, and $\gamma_2$ bounds a disc in $\Omega'$ containing no point of $A$. Twists about such curves are trivial in $\PMod(S^2,A)$, so $[\tilde T]=1$, i.e.\ $f^*\circ T_\gamma=f^*$.
\end{remark}

\section{Transport to an arbitrary \texorpdfstring{$f$}{f}: proof of Theorem~\ref{thm:main}(b)}

Fix $t\in W$ and a path $t_s$ ($s\in[0,1]$) in $W$ from $t_0=\tfrac14$ to $t$. Write $f_s=f_{t_s}$, $B_s=\{0,\infty,1,t_s,v_s\}$ with $v_s=v(t_s)$, and $A_s=B_s\setminus\{v_s\}$.

\begin{lemma}[isotopy]\label{lem:isotopy}
There is a jointly continuous isotopy $h_s$ of $\Ch$ with $h_0=\mathrm{id}$, fixing $0,1,\infty$, with $h_s(t_0)=t_s$ and $h_s(v_0)=v_s$.
\end{lemma}

\begin{proof}
For $|q|\le\tfrac12$ the map $P_q(z)=z+(1-|z|)q$ is a homeomorphism of the closed unit disc that is the identity on the boundary and sends $0\mapsto q$. It maps the disc into itself, since $|P_q(z)|\le|z|+(1-|z|)|q|$. It is injective, since $P_q(z)=P_q(z')$ gives $|z-z'|=\bigl||z|-|z'|\bigr|\,|q|\le|z-z'||q|$. It is surjective by invariance of domain, the boundary being fixed. Finally $(q,z)\mapsto(q,P_q(z))$ is a continuous bijection of a compact Hausdorff space, hence a homeomorphism.

Let $\delta>0$ be less than one third of the minimum, over $s\in[0,1]$, of the mutual distances among $0,1,t_s,v_s$ (positive by compactness, since $t_s\in W$ keeps these four points distinct and finite). By uniform continuity choose $0=s_0<s_1<\dots<s_n=1$ with $|t_s-t_{s_k}|<\delta/2$ and $|v_s-v_{s_k}|<\delta/2$ for $s\in[s_k,s_{k+1}]$. On $[s_k,s_{k+1}]$ let $D_t,D_v$ be the closed discs of radius $\delta$ about $t_{s_k},v_{s_k}$; they are disjoint and avoid $0,1,\infty$. Let $\Pi^t_s$ be the affine conjugate of $P_q$ to $D_t$ with $q=(t_s-t_{s_k})/\delta$, extended by the identity off $D_t$ (so it moves $t_{s_k}$ to $t_s$), and similarly $\Pi^v_s$ on $D_v$. Set $h_s=\Pi^t_s\Pi^v_s h_{s_k}$. At $s=s_k$ we have $q=0$ and $P_0=\mathrm{id}$, so the definitions agree at the junctions; joint continuity follows from that of $(q,z)\mapsto P_q(z)$.
\end{proof}

\begin{lemma}[lift]\label{lem:liftiso}
There is a continuous family of homeomorphisms $\tilde h_s$ with $f_s\tilde h_s=h_sf_0$, $\tilde h_0=\mathrm{id}$, and $\tilde h_s=h_s$ on $A_0$.
\end{lemma}

\begin{proof}
Write $f_s=M_sQN_s$ with $Q(z)=z^2$, $M_s=M_{t_s}$, $N_s=N_{t_s}$, by \eqref{eq:factor}. Put $k_s=M_s^{-1}h_sM_0$. Then $k_s$ fixes $0$ and $\infty$, since $M_0(0)=\infty$, $M_0(\infty)=v_0$, $h_s(v_0)=v_s$ and $M_s(\infty)=v_s$. The map $K(s,z)=k_s(z^2)$ on $[0,1]\times\mathbb C^*$ lifts through the covering $Q\colon\mathbb C^*\to\mathbb C^*$, because the domain retracts to $s=0$, where $K=Q$. This gives $\tilde k_s$ with $Q\tilde k_s=k_sQ$ and $\tilde k_0=\mathrm{id}$. By uniqueness of lifts $\tilde k_s$ is odd ($z\mapsto-\tilde k_s(-z)$ is a lift agreeing with $\tilde k_s$ at $s=0$), hence injective: $\tilde k_s(z)=\tilde k_s(z')$ forces $z'=\pm z$, and $z'=-z$ would give $\tilde k_s(z)=0$. It is surjective, again by oddness, and it extends by $0\mapsto0$, $\infty\mapsto\infty$ to a homeomorphism of $\Ch$. Put $\tilde h_s=N_s^{-1}\tilde k_sN_0$. Then $f_s\tilde h_s=M_sQ\tilde k_sN_0=M_sk_sQN_0=h_sf_0$.

On $A_0$: $\tilde h_s(0)$ is a preimage under $f_s$ of $h_s(a_1)=\infty$, whose unique preimage is $0$. For $i=1,2,3$, $\tilde h_s(a_i)$ is a continuous selection from the two distinct preimages under $f_s$ of $a_{i+1}(s)=h_s(a_{i+1})$, which is not a critical value. These preimages move continuously, one of them is $a_i(s)$, and the selection starts at $a_i$. So it equals $a_i(s)=h_s(a_i)$ for all $s$.
\end{proof}

Therefore $\kappa_s=h_s^{-1}\tilde h_s$ is an isotopy rel $A_0$ starting at the identity, so $\tilde h_1\simeq h_1$ rel $A_0$.

\begin{proposition}\label{prop:transport}
Let $h=h_1$ and $\tilde h=\tilde h_1$, and define $h^*[\varphi]=[\varphi h]$ from $\Teich(B_1)$ to $\Teich(B_0)$, and likewise from $\Teich(A_1)$ to $\Teich(A_0)$. Then:
\begin{enumerate}[label=\textup{(\roman*)},nosep]
\item $i_*h^*=h^*i_*$ and $f_0^*h^*=h^*f_t^*$;
\item hence $h^*(\D_{f_t})=\D_{f_0}$, with components mapping to components;
\item $h^*\circ T_{h(\gamma)}=T_\gamma\circ h^*$.
\end{enumerate}
\end{proposition}

\begin{proof}
The first identity in (i) is clear. If $\varphi f_t=F\psi$, then $\varphi hf_0=\varphi f_t\tilde h=F\psi\tilde h$, so $f_0^*[\varphi h]=[\psi\tilde h]=[\psi h]$ because $\tilde h\simeq h$ rel $A_0$. For (ii), $h^*$ is a homeomorphism with inverse $(h^{-1})^*$. For (iii), $h^*(hgh^{-1}\cdot[\varphi])=[\varphi hg^{-1}]=g\cdot h^*[\varphi]$, and $hT_\gamma h^{-1}=T_{h(\gamma)}$.
\end{proof}

Hence $T_{h_1(\gamma)}$ preserves $\D_{f_t}$, and every orbit of $\langle T_{h_1(\gamma)}\rangle$ on its components is infinite, by Theorem~\ref{thm:main}(a).

For an abstract $f$, take $M$ from Proposition~\ref{prop:normal}; it maps $B$ to $B_t$, since $M(v)=M(f(c_2))=f_t(M(c_2))=v(t)$. The map $M^*\colon[\varphi]\mapsto[\varphi M]$ satisfies the same identities as in Proposition~\ref{prop:transport}, with $h=\tilde h=M$, because $f_tM=Mf$. So $\gamma_f=M^{-1}h_1(\gamma)$ has the properties stated in Theorem~\ref{thm:main}(b).

\paragraph{Real $t\in I=\bigl(0,\tfrac{3-\sqrt5}2\bigr)$.} The interval $I$ lies in $W$, and on $I$ we have $v(t)<0$ (the numerator of $-v$ is a nonzero square and $4t(1-t)(1-2t)>0$). So the order $v_s<0<t_s<1<\infty$ persists along the segment from $\tfrac14$ to $t$. Instead of the isotopy of Lemma~\ref{lem:isotopy} take $h_s(x+iy)=\phi_s(x)+iy$, where $\phi_s$ is the increasing piecewise-linear map of $\mathbb R$ fixing $0$, $1$ and everything outside a large interval, with $t_0\mapsto t_s$ and $v_0\mapsto v_s$; it is jointly continuous and has all the properties used in Lemma~\ref{lem:liftiso} and Proposition~\ref{prop:transport}. It preserves $\mathbb R$, both half-planes, and the intervals between marked points. So $h_1(\alpha)$ has the same crossing pattern, with $\tfrac14$ replaced by $t$ and $v$ by $v(t)$. Pulling back by $M$ gives the description in Theorem~\ref{thm:main}(b). For non-real $t$ the curve is given as $M^{-1}h_1(\gamma)$ with $h_1$ the explicit isotopy of Lemma~\ref{lem:isotopy} along a chosen path in $W$; every choice of path works. \qed

\begin{thebibliography}{99}

\bibitem{Ahlfors} L.~V.~Ahlfors, \emph{Complex Analysis}, 3rd ed., McGraw-Hill, New York, 1979.

\bibitem{Bers} L.~Bers, Fiber spaces over Teichm\"uller spaces, \emph{Acta Math.} \textbf{130} (1973), 89--126.

\bibitem{Birman} J.~S.~Birman, Mapping class groups and their relationship to braid groups, \emph{Comm. Pure Appl. Math.} \textbf{22} (1969), 213--238.

\bibitem{BEKP} X.~Buff, A.~Epstein, S.~Koch and K.~Pilgrim, On Thurston's pullback map, in: \emph{Complex Dynamics: Families and Friends} (D.~Schleicher, ed.), A K Peters, Wellesley, MA, 2009.

\bibitem{Chandra} K.~Chandrasekharan, \emph{Elliptic Functions}, Grundlehren der mathematischen Wissenschaften 281, Springer, Berlin, 1985.

\bibitem{DH} A.~Douady and J.~H.~Hubbard, A proof of Thurston's topological characterization of rational functions, \emph{Acta Math.} \textbf{171} (1993), 263--297.

\bibitem{FM} B.~Farb and D.~Margalit, \emph{A Primer on Mapping Class Groups}, Princeton University Press, Princeton, NJ, 2012.

\bibitem{HK} E.~Hironaka and S.~Koch, A disconnected deformation space of rational maps, \emph{J. Mod. Dyn.} (2017).

\bibitem{Hubbard} J.~H.~Hubbard, \emph{Teichm\"uller Theory and Applications to Geometry, Topology, and Dynamics, Vol.~1: Teichm\"uller Theory}, Matrix Editions, Ithaca, NY, 2006.

\end{thebibliography}

\end{document}
